\documentclass[11pt]{article}

\usepackage[margin=1in]{geometry}
\usepackage{microtype}
\usepackage{amsmath}
\usepackage{graphicx}
\usepackage{booktabs}
\usepackage[hidelinks]{hyperref}

\graphicspath{{figures/}}

\newcommand{\natex}{\texttt{natex}}

\title{Bunching Below the Line: Missing Mass Above the EU AI Act's
$10^{25}$-FLOP Threshold}
\author{Hauke Hillebrandt\thanks{University College London. Email:
\texttt{ucjthhi@ucl.ac.uk}. This paper and the underlying \natex{} software were
prepared with substantial assistance from Anthropic's Claude models; the author
reviewed the analyses and text and is responsible for all remaining errors.}}
\date{July 2026 \\ \small analysis pass of record: 2026-07, fully
deterministic (no RNG)}

\begin{document}

\maketitle

\begin{abstract}
\noindent After the EU AI Act entered into force on 1 August 2024, the density
of new AI models just \emph{above} the Act's $10^{25}$-FLOP systemic-risk
presumption line collapses relative to just below it. In Epoch AI's compute
panel (719 models with training-compute estimates, January 2023 to June 2026),
a McCrary-style binned-Poisson density test at the statutory cutoff
$\log_{10}\mathrm{FLOP}=25.0$ gives $\theta=-1.334$ (SE $0.878$, $p=0.129$) in
the primary $\pm 0.5$-dex window and is negative in all 15 window~$\times$~bin
specifications (10 of 15 with $p<0.05$); the pre-Act period is never
significant (0 of 15). A Fisher exact below/above~$\times$~pre/post contrast
puts the shift at odds ratio $0.173$ (95\% CI $[0.048, 0.619]$, $p=0.00715$):
at the pre-Act above:below ratio, $28.9$ post-Act models are expected in
$(25.0, 25.5]$ and 5 are observed --- an $82.7\%$ deficit. The deficit is
specific to the statutory line (placebo thresholds at $24.5$ and $25.5$ are
null or opposite-signed, with mass piling up just below $25.0$), specific to
the post-Act period (a placebo split date inside the pre-Act sample gives
OR $0.90$, $p=1.0$), and correctly timed: the above-line count collapses from
2025H1, after entry into force and before the August 2025 applicability date
for general-purpose AI obligations. Two results are nulls and are stated as
nulls: the preferred narrow-window pre/post interaction in the discontinuity
is $p=0.729$, significant only at the widest bandwidth, and the pre-Act
density test never rejects. With only 5--13 post-Act models above the line,
and an Act-induced \emph{disclosure} response that would mimic bunching
exactly, we grade this as moderate evidence of bunching --- behavioral or
reporting --- at the EU AI Act's compute threshold, not proof of a
training-compute response.
\end{abstract}

\section{Introduction}\label{sec:intro}

Article 51(2) of the EU AI Act \cite{euaiact2024} presumes that a
general-purpose AI model poses ``systemic risk'' when the cumulative compute
used for its training exceeds $10^{25}$ floating-point operations, attaching
notification, evaluation, and risk-mitigation obligations to models above the
line. The Act entered into force on 1 August 2024 and its general-purpose AI
obligations became applicable on 2 August 2025. A bright-line quantity
threshold of this kind is, in the language of public economics, a
\emph{notch}: crossing it discretely changes the regulatory burden, so agents
with values just above the cutoff have an incentive to locate just below it.

A large empirical literature shows that agents do exactly this.
\cite{saez2010} developed bunching estimation at kinks in the US income-tax
schedule; \cite{klevenwaseem2013} extended it to notches, showing that a
discrete jump in liability produces excess mass just below the cutoff and
\emph{missing mass} just above it; \cite{kleven2016} surveys the method and
its spread beyond taxation to regulation and private-sector prices. The
econometric diagnostic for such sorting at a cutoff is the density
discontinuity test of \cite{mccrary2008}: if agents manipulate the running
variable, its density jumps at the threshold.

Whether training compute is a good handle for AI governance is actively
debated. \cite{sastry2024} argue that compute is uniquely governable ---
detectable, quantifiable, and produced by a concentrated supply chain ---
and \cite{heimkoessler2024} defend training-compute thresholds, including the
EU's $10^{25}$-FLOP line, as the most workable trigger for regulatory
oversight. \cite{hooker2024} counters that compute thresholds are shortsighted
and invite exactly the strategic responses studied in the bunching
literature. On the empirical side, \cite{kumarmanning2025} project model
counts above the EU threshold \emph{forward} from pre-Act scaling trends,
forecasting 103--306 models above $10^{25}$ FLOP by 2028 --- a superlinear
increase. To our knowledge, no published estimate tests whether the observed
post-Act density of models near the line already departs from that trend.
This paper supplies that test: a McCrary-style density discontinuity estimate
at the statutory cutoff, before versus after the Act's entry into force, with
threshold, date, and timing placebos.

\section{Data}\label{sec:data}

The input panel is built from Epoch AI's \emph{Data on AI Models} database
\cite{epochai}, keeping every model with a non-null positive training-compute
estimate and a release date from 2 January 2023 to 4 June 2026: 719 models
with columns model, date, and $\log_{10}$ training FLOP. Event dates are the
Act's entry into force (2024-08-01) and the applicability of general-purpose
AI obligations (2025-08-02). The pre-Act period is 2023-01 to 2024-07
($n=21$ models within $\pm 0.5$ dex of the cutoff) and the post-Act period is
2024-08 onward ($n=52$ in-window). Compute values for recent frontier models
are largely Epoch estimates rather than disclosures, and models without any
compute estimate ($\approx 2{,}800$ of $\approx 3{,}500$ database rows) are
necessarily excluded --- a selection channel we return to in
Section~\ref{sec:conclusion}. There is no value heaping near the cutoff: at
most 3 duplicate $\log_{10}$-compute values occur within $[24.5, 25.5]$, so
the results below are not an artifact of rounded compute estimates.

\section{Design and Methods}\label{sec:methods}

\paragraph{Density discontinuity at the statutory line.} The main estimate is
the \natex{} toolkit's \cite{natex} McCrary-style \cite{mccrary2008}
binned-Poisson density test, run directly at the statutory cutoff: model
counts in equal-width bins of the signed distance
$\log_{10}\mathrm{FLOP}-25.0$ are fit by a Poisson regression with a linear
trend in distance and a jump term at zero; $\theta$ is the jump in log density
at the cutoff, so $\theta<0$ means a deficit just above the line ---
avoidance bunching below it. The primary specification uses a $\pm 0.5$-dex
window and 10 bins; sensitivity covers the full grid $h\in\{0.5,0.75,1.0\}
\times \text{bins}\in\{8,10,12,16,20\}$ (15 specifications per period).
Because the cutoff is fixed by statute rather than searched for, the
frozen-geometry caveat that applies to \natex{}'s post-discovery density
checks does not bite here.

\paragraph{Contrasts and placebos.} The pre/post change in bunching is
measured two ways: a Fisher exact test on the $2\times 2$ table of
below/above counts in the $\pm 0.5$-dex window by period, and a
$\theta$-difference (interaction) $z$-test between the period-specific
density jumps. Specificity is probed with placebo thresholds at $24.5$ and
$25.5$ and a placebo split date of 2023-08-01, which places both
``periods'' inside the pre-Act sample.

\paragraph{Honest inference.} Standard errors for $\theta$ are Wald standard
errors from the binned Poisson fit and can be optimistic when the underlying
density is strongly curved within the window; at the widest bandwidth
($h=1.0$) curvature of the steeply falling compute density can leak into
$\theta$, which is why the narrow-window specification is primary and the
wide-window significance is discounted below. All counts, tests, and figures
derive deterministically from the input panel: there is no random number
generation anywhere in the pipeline.

\section{Results}\label{sec:results}

\begin{figure}[t]
\centering
\includegraphics[width=\textwidth]{fig1}
\caption{(a) Model counts in 0.25-dex bins of $\log_{10}$ training compute,
pre-Act (top) versus post-Act (bottom), colored by side of the statutory
$10^{25}$-FLOP line (dashed). Post-Act, mass accumulates just below the line
while the first half-dex above holds 5 models against $28.9$ expected at the
pre-Act ratio; the frontier jumps far above the notch. (b) Below/above counts
per half-year in the $\pm 0.5$-dex window: the above-line count (labeled)
collapses from 2025H1, after entry into force and before GPAI applicability,
while the just-below count keeps growing.}
\label{fig:fig1}
\end{figure}

\paragraph{A post-Act deficit just above the line.} Post-Act (2024-08-01
onward), the primary density test at $\log_{10}=25.0$ gives
$\theta=-1.334$ (SE $0.878$, $p=0.129$) with $n=52$ in-window models split
47 below / 5 above. Across the full sensitivity grid $\theta$ is negative in
15 of 15 specifications (range $-1.249$ to $-2.566$) and significant at the
5\% level in 10 of 15 --- all $h=0.75$ and $h=1.0$ specifications, with
significant $p$-values from $0.0010$ to $0.0143$. The pre-Act period shows no
such pattern: $\theta=-0.880$ (SE $0.973$, $p=0.366$), $n=21$ split 13/8,
and 0 of 15 sensitivity specifications reject ($\theta$ from $-1.173$ to
$+0.494$, minimum $p=0.243$).

\paragraph{The pre/post contrast.} The Fisher exact test on below/above
$\times$ pre/post counts in the $\pm 0.5$-dex window --- table
$[[13, 8], [47, 5]]$ --- gives OR $=0.173$ (95\% CI $[0.048, 0.619]$,
$p=0.00715$). At the pre-Act above:below ratio, $28.9$ post-Act models are
expected in $(25.0, 25.5]$; 5 are observed, so roughly 24 models are
``missing'' --- an $82.7\%$ deficit. The formal interaction between the
period-specific density jumps is a null at the primary specification and is
stated as one: $\Delta\theta=-0.454$ (SE $1.311$, $z=-0.346$, $p=0.729$).
Across the grid the interaction is always negative at $h\geq 0.75$ but
significant only at $h=1.0$ with 12/16/20 bins ($p=0.0249$/$0.0049$/$0.0342$;
$h=1.0$ with 10 bins gives $p=0.0639$), where curvature leakage is a concern
(Section~\ref{sec:methods}). The pre-period, with only 21--47 in-window
models, is too small to pin down the counterfactual discontinuity.

\paragraph{Placebos.} The deficit is specific to the statutory line and
period. At placebo threshold $24.5$ the same pre/post Fisher contrast gives
OR $=1.637$, $p=0.248$; restricted to the post-GPAI period it gives
OR $=2.510$, $p=0.0386$ with the \emph{opposite} sign --- more mass in
$[24.5, 25.0)$ than $[24.0, 24.5)$, i.e.\ mass piling up just below $25.0$,
consistent with bunching at the line and inconsistent with a generic
high-compute slowdown. At placebo threshold $25.5$: OR $=4.00$, $p=0.350$,
direction opposite to the $25.0$ effect. The placebo split date 2023-08-01,
with both halves pre-Act, gives OR $=0.90$, $p=1.0$ at the true threshold ---
no pseudo-effect.

\paragraph{Timing.} Below:above counts per half-year in the $\pm 0.5$-dex
window are 2023H1 2:2, 2023H2 3:1, 2024H1 6:4, 2024H2 11:3, then 2025H1
17:1, 2025H2 15:1, 2026H1 6:1 (Figure~\ref{fig:fig1}b). The collapse begins
in 2025H1 --- after entry into force (August 2024), before the obligations
became applicable (August 2025) --- consistent with anticipatory compliance:
the $10^{25}$ presumption was public from July 2024 and attaches to models
placed on the market from August 2025. By year, $[25.0, 25.5)$ holds 3, 7,
2, 1 models in 2023--2026 while $[24.5, 25.0)$ grows 5, 17, 32, 6. Secular
compute growth biases \emph{against} the finding --- later cohorts should put
more, not less, mass above any fixed line, and do at $24.5$. Meanwhile 13
post-Act models sit at or above $10^{25}$ at some height, with the frontier
jumping far above the notch (Grok~3 at $\log_{10}=26.54$, GPT-4.5 at
$26.58$, Llama~4 Behemoth preview at $25.71$, Grok~4 at $26.70$, GPT-5 at
$25.82$) --- the classic bunching signature in which a notch distorts only
marginal decisions \cite{klevenwaseem2013}.

\begin{table}[t]
\centering
\small
\caption{Headline estimates. Density tests: binned-Poisson jump $\theta$ in
log density at the statutory cutoff $\log_{10}\mathrm{FLOP}=25.0$; the
primary window is $\pm 0.5$ dex with 10 bins, and the sensitivity grid
crosses $h\in\{0.5,0.75,1.0\}$ with $\{8,10,12,16,20\}$ bins. Fisher tests
use below/above counts in the $\pm 0.5$-dex window by period.}
\label{tab:main}
\begin{tabular}{lrrrl}
\toprule
Quantity & Estimate & SE / CI & $p$ & notes \\
\midrule
\multicolumn{5}{l}{\emph{Panel A: density jump at the statutory line}}\\
$\theta$, pre-Act (2023-01..2024-07)  & $-0.880$ & $0.973$ & $0.366$ & $n=21$, 13/8; 0/15 specs $p<0.05$ \\
$\theta$, post-Act (2024-08..)        & $-1.334$ & $0.878$ & $0.129$ & $n=52$, 47/5; 10/15 specs $p<0.05$ \\
$\Delta\theta$ (post $-$ pre)         & $-0.454$ & $1.311$ & $0.729$ & \emph{null at primary spec} \\
\midrule
\multicolumn{5}{l}{\emph{Panel B: Fisher exact below/above $\times$ pre/post}}\\
threshold $25.0$, Act date       & OR $0.173$ & $[0.048, 0.619]$ & $0.00715$ & $[[13,8],[47,5]]$; $-82.7\%$ \\
placebo threshold $24.5$         & OR $1.637$ & ---     & $0.248$  & post-GPAI: OR $2.510$, $p=0.0386$ \\
placebo threshold $25.5$         & OR $4.00$  & ---     & $0.350$  & opposite direction \\
placebo date 2023-08-01          & OR $0.90$  & ---     & $1.0$    & both halves pre-Act; no effect \\
\bottomrule
\end{tabular}

\medskip
\parbox{0.92\textwidth}{\footnotesize Notes: $\theta<0$ is a density deficit
just above the cutoff. SEs are Wald standard errors from the binned Poisson
fit. Significant sensitivity $p$-values for the post-Act $\theta$ range from
$0.0010$ to $0.0143$ (all $h=0.75$ and $h=1.0$ specifications); the
interaction is significant only at $h=1.0$ with $\geq 12$ bins
($p=0.0049$--$0.0342$). Placebo-threshold odds ratios use the same
$\pm 0.5$-dex construction centered on the placebo line.}
\end{table}

\section{Caveats and Conclusion}\label{sec:conclusion}

Three caveats bound the claims. First, \emph{small n}: only 5 models sit in
$(25.0, 25.5]$ and 13 at any height above $10^{25}$ post-Act; the preferred
narrow-window pre/post interaction is $p=0.729$, and the interaction is
significant only at $h=1.0$ with $\geq 12$ bins ($p$ from $0.0049$ to
$0.0342$), where curvature of the steeply falling density can leak into
$\theta$. Second, \emph{measurement, not behavior}: post-2024 frontier
compute values are largely Epoch estimates, and models without a compute
estimate are excluded (719 of $\approx 3{,}500$ rows have compute). If the
Act suppresses \emph{disclosure} of parameters and token counts near the
line, the identical density gap appears with no training-compute response at
all; distinguishing a compliance response from a reporting response requires
disclosure-side data and is the designated follow-up. Third, a
\emph{confounded trend break}: the deficit's onset (2025H1) coincides with
the industry shift toward inference-time scaling and distilled or
mixture-of-experts models that followed the late-2024 o1-style reasoning
models, which independently reduced demand for $10^{25}$-plus pre-training
runs among non-frontier labs --- though this confound predicts a broad
high-compute slowdown, not a deficit specific to the statutory line with
mass accumulating immediately below it.

The conclusion is graded accordingly. The post-Act deficit just above the EU
AI Act's $10^{25}$-FLOP line is sign-consistent in all 15 density
specifications, large (Fisher $p=0.00715$, $82.7\%$ missing mass), specific
to the statutory threshold and period, and correctly timed, with the
marginal-crossers-vanish, frontier-jumps-far-above signature that notch
responses produce. But with 5--13 post-Act models above the line, a null
preferred narrow-window interaction, and an Act-induced disclosure response
that would mimic bunching exactly, it is moderate evidence of bunching ---
behavioral or reporting --- rather than proof of a training-compute
response. Either reading is informative for the compute-threshold debate
\cite{heimkoessler2024,hooker2024}: the forward projections of
\cite{kumarmanning2025} assumed threshold-blind scaling, and the data just
above the line already look threshold-aware.

\paragraph{Reproducibility.} All estimates derive deterministically from the
public Epoch AI panel; no dataset files are committed.
Figure~\ref{fig:fig1} regenerates from the committed
\texttt{figures/make\_fig.py}, which asserts the headline estimates against
the numbers of record before drawing.

\footnotesize
\begin{thebibliography}{9}
\setlength{\itemsep}{1pt}\setlength{\parskip}{0pt}

\bibitem{epochai}
Epoch AI (2026).
Data on AI models.
\url{https://epoch.ai/data/ai-models} (compute panel extracted June 2026).

\bibitem{euaiact2024}
European Union (2024).
Regulation (EU) 2024/1689 of the European Parliament and of the Council
laying down harmonised rules on artificial intelligence (Artificial
Intelligence Act).
\emph{Official Journal of the European Union}, L series, 12 July 2024.

\bibitem{heimkoessler2024}
Heim, L., and Koessler, L. (2024).
Training compute thresholds: Features and functions in AI regulation.
arXiv:2405.10799.

\bibitem{natex}
Hillebrandt, H. (2026).
\emph{natex: automated natural-experiment discovery and estimation}.
Software. \url{https://github.com/HaukeHillebrandt/natex}

\bibitem{hooker2024}
Hooker, S. (2024).
On the limitations of compute thresholds as a governance strategy.
arXiv:2407.05694.

\bibitem{kleven2016}
Kleven, H.~J. (2016).
Bunching.
\emph{Annual Review of Economics}, 8, 435--464.

\bibitem{klevenwaseem2013}
Kleven, H.~J., and Waseem, M. (2013).
Using notches to uncover optimization frictions and structural elasticities:
Theory and evidence from Pakistan.
\emph{Quarterly Journal of Economics}, 128(2), 669--723.

\bibitem{kumarmanning2025}
Kumar, I., and Manning, S. (2025).
Trends in frontier AI model count: A forecast to 2028.
arXiv:2504.16138.

\bibitem{mccrary2008}
McCrary, J. (2008).
Manipulation of the running variable in the regression discontinuity design:
A density test.
\emph{Journal of Econometrics}, 142(2), 698--714.

\bibitem{saez2010}
Saez, E. (2010).
Do taxpayers bunch at kink points?
\emph{American Economic Journal: Economic Policy}, 2(3), 180--212.

\bibitem{sastry2024}
Sastry, G., Heim, L., Belfield, H., Anderljung, M., Brundage, M., Hazell, J.,
O'Keefe, C., Hadfield, G.~K., et al. (2024).
Computing power and the governance of artificial intelligence.
arXiv:2402.08797.

\end{thebibliography}

\end{document}
